\section*{अध्याय \ref{ch:b3:conservation-biology} --- संरक्षण जीवविज्ञान}

\begin{solution}{exo:b3:conservation-biology:1}
\omterm{def:b3:conservation-biology:small}{जनसांख्यिकीय प्रसंभाव्यता}: थोड़े-से व्यष्टियों के जन्मों, मृत्युओं तथा
लिंगों में संयोग --- इक्यावन पर काकापो, जिसमें मादाओं से अधिक नर थे।
\omterm{def:b3:conservation-biology:small}{पर्यावरणीय प्रसंभाव्यता}: कोई बुरा वर्ष जो हर व्यष्टि पर एक साथ पड़े ---
कोई सूखा, कोई रोग, कुछ सौ की किसी समष्टि पर पड़ती कोई कड़ी सर्दी। \omterm{def:b3:conservation-biology:small}{ऐली
प्रभाव}: कम घनत्व पर गिरती प्रति-व्यष्टि वृद्धि --- यात्री कबूतर, जिसकी
बस्तियाँ विरल हो जाने पर प्रजनन न कर सकीं, और जिसके बचे हुए परभक्षी अब
अभिभूत नहीं होते थे।
\end{solution}

\begin{solution}{exo:b3:conservation-biology:2}
किसी आदर्श समष्टि का वह आकार --- अचल, यादृच्छिक संगम वाली, बराबर लिंगों
वाली, प्वासों कुल-आकारों वाली --- जो असली समष्टि की ही दर से
विषमयुग्मजता खोती। प्रजननशील नरों तथा मादाओं की असमान संख्याएँ; आकार में
उतार-चढ़ाव, जो सबसे छोटी पीढ़ियों को भार देते हैं; प्वासों से ऊपर
कुल-आकार का प्रसरण, जिसमें कुछ ही व्यष्टि अधिकांश बच्चे देते हैं (साथ ही
अयादृच्छिक संगम तथा अतिव्यापी पीढ़ियाँ)।
\end{solution}

\begin{solution}{exo:b3:conservation-biology:3}
$S = cA^{z}$। रखा गया अंश $= (A'/A)^{z}$: $0.5^{0.25} = 0.84$;
$0.1^{0.25} = 0.56$; $0.01^{0.25} = 0.32$ --- क्षेत्रफल का सौवाँ भाग भी
अंततः जातियों का तिहाई रखता है।
\end{solution}

\begin{solution}{exo:b3:conservation-biology:4}
$N_{e} \ge 50$ अंतःप्रजनन की वृद्धि को प्रति पीढ़ी $1/100$ से नीचे रखता
है, यानी वह दर जिस पर पशु-प्रजनक अवसाद को सहनीय पाते हैं: वह सुयोग्यता की
अल्पकालिक हानि से बचाता है। $N_{e} \ge
500$ अपवाह से विविधता की हानि को उत्परिवर्तन से उसकी पूर्ति के सामने
संतुलित करता है, ताकि समष्टि वह मात्रात्मक विविधता रखे जो उसे अनुकूलन के
लिए चाहिए: वह दीर्घ काल की रक्षा करता है। असली समष्टियों को गणना में इन
संख्याओं से पाँच से दस गुना चाहिए।
\end{solution}

\begin{solution}{exo:b3:conservation-biology:5}
$N_{e} = 4\times 8\times 120/128 = 30$; विषमयुग्मजता प्रति पीढ़ी $1/60 =
\qty{1.7}{\%}$ गिरती है; चौथाई तब खोया जब $(1 - 1/60)^{t} =
0.75$, यानी $t = \ln 0.75/\ln(59/60) = 17$ पीढ़ियाँ।
\end{solution}

\begin{solution}{exo:b3:conservation-biology:6}
$1/N_{e} = \tfrac{1}{5}(4/1000 + 1/20) = 0.0108$, $N_{e} = 93$: बीस पर
बिताई एक ही पीढ़ी ने पाँच पीढ़ियों को तिरानबे जैसा बरतवा दिया। बची हुई
विषमयुग्मजता $(1 - 1/186)^{5} = 0.973$, जबकि अचल हज़ार के लिए $(1 -
1/2000)^{5} = 0.9975$।
\end{solution}

\begin{solution}{exo:b3:conservation-biology:7}
$F = 1 - (1 - 1/20)^{5} = 0.23$; सुयोग्यता $5\times 2.3 = \qty{11}{\%}$
नीचे। किसी टेक्सस मादा और किसी फ़्लोरिडा नर की हर संतान का $F = 0$ होता
है; यदि वे आठ नवागंतुक प्रजननशील मादाओं के दो-पाँचवें भाग हों, तो अगली
पीढ़ी का दो-पाँचवाँ भाग बहिःप्रजनित होता है और माध्य $F$ एक ही पीढ़ी में
गिरकर लगभग $0.6\times 0.23 = 0.14$ रह जाता है --- और वापसी ने यही
दिखाया।
\end{solution}

\begin{solution}{exo:b3:conservation-biology:8}
$r = \ln(100/31)/8 = \qty{0.15}{yr^{-1}}$। 2020 तक चलती रहती:
$100\,\mathrm{e}^{0.15\times 17} \approx 1200$। उद्यान ने लगभग सौ थामे:
क्षेत्र भर गए, एल्क घटे, झुंडों ने एक-दूसरे के सदस्य मारे, और खुजली तथा
डिस्टेंपर फैले --- ख़ाली विस्तार भर जाने के बाद घनत्व-निर्भरता बैठ गई।
\end{solution}

\begin{solution}{exo:b3:conservation-biology:9}
(a) दस वर्षों में \qty{60}{\%} का ह्रास \qty{50}{\%} से अधिक है:
संकटग्रस्त। (b) $180 < 250$ परिपक्व व्यष्टि: अतिसंकटग्रस्त। (c) विस्तार
\qty{3000}{km^2}, \qty{5000}{} से नीचे: संकटग्रस्त। (d) \num{30000}
व्यष्टि और \qty{10}{\%} का ह्रास: कोई भी देहली पार नहीं हुई --- संकट में
नहीं (अधिक से अधिक संकट के निकट, जिस पर नज़र रखी जाए)।
\end{solution}

\begin{solution}{exo:b3:conservation-biology:10}
सारे $N$ युग्म $(1/2)^{N}$ प्रायिकता से चूकते हैं: $0.25$, $0.031$,
$0.001$, $10^{-6}$। जनसांख्यिकीय जोखिम $N$ के साथ चरघातांकी रूप से गिरता
है और कुछ दर्जन व्यष्टियों से आगे नगण्य है; उससे ऊपर जो जोखिम बचता है वह
पर्यावरणीय है, जो इतनी तेज़ी से नहीं गिरता, क्योंकि वह सारे व्यष्टियों
पर एक साथ पड़ता है।
\end{solution}

\begin{solution}{exo:b3:conservation-biology:11}
दसवें भाग का कोई एक खंड: मूल जातियों का $0.1^{0.25} = 0.56$। दस खंड
$0.56$ (यदि सब एक-सी जातियाँ थामें) और, सिद्धांततः, पूरे वन से अधिक (यदि
हर एक भिन्न जातियाँ थामे --- जो क्षेत्रीय कोश से आगे असंभव है) के बीच
थामते हैं। वे कहाँ गिरते हैं यह खंडों के बीच के आवर्तन पर निर्भर करता
है: भिन्न खंडों में भिन्न आवास कुल को ऊपर धकेलते हैं; जिन जातियों को बड़े
क्षेत्रफल या भीतरी आवास चाहिए, या जो उन्हीं \omterm{prop:b3:conservation-biology:area}{किनारा प्रभावों} से हर खंड से
खो जाती हैं, वे उसे नीचे धकेलती हैं; और हर खंड का विलोपन-ऋण अलग से
चुकाया जाता है, इसलिए दीर्घ काल का कुल प्रायः पूरे वन के कुल से नीचे
रहता है।
\end{solution}

\begin{solution}{exo:b3:conservation-biology:12}
अपवाह विषमयुग्मजता को प्रति पीढ़ी $1/2N_{e}$ दर से हटाता है; आप्रवासी
जीन-कोश के $m$ अंश को कहीं और से खींचे गए युग्मविकल्पियों से बदल देते
हैं। $mN_{e} \approx 1$ के साथ, $m \approx 1/N_{e}$, यानी हानि की दर से
दुगुना, इसलिए विविधता उसके क्षय से तेज़ी से भरती रहती है और समष्टि
युग्मविकल्पी-बारंबारता में स्रोत के पास बनी रहती है ($F_{ST} = 1/(1
+ 4N_{e}m) \approx 0.2$)। लागत: स्थानीय रूप से पक्ष पाए युग्मविकल्पी
प्रति पीढ़ी $m$ से पतले होते हैं, इसलिए स्थानीय अनुकूलन केवल उन लक्षणों
के लिए बचता है जिनका वरण गुणांक लगभग $1/N_{e}$ से अधिक हो; दुर्बल वरण
वाले स्थानीय भेद बहा दिए जाते हैं।
\end{solution}

\begin{solution}{pb:b3:conservation-biology:1}
\textbf{1.} $N_{e} = 4\times 20\times 40/60 = 53$, जबकि गणना 60 की है।
\textbf{2.} $1/N_{e} = \tfrac{1}{5}(1/500 + 1/100 + 1/60 + 1/51 + 1/60) =
0.0130$, $N_{e} = 77$; 51 तथा 60 वाली अड़चन-पीढ़ियाँ हावी हैं, वह 500
शायद ही गिनती में आता है।
\textbf{3.} $F = 1 - (1 - 1/154)^{5} = 0.032$; रैखिक $5/154 = 0.032$।
\textbf{4.} $6\times 0.32 = \qty{1.9}{\%}$।
\textbf{5.} $(1 - 1/154)^{5} = 0.968$: \qty{96.8}{\%} बचा।
\textbf{6.} $N_{e} = 53$ पर दस पीढ़ियाँ: वृद्धि $1 - (1 -
1/107)^{10} = 0.090$, कुल $F = 1 - 0.968\times 0.910 = 0.12$;
सुयोग्यता-हानि $6\times 1.2 = \qty{7}{\%}$; यानी कोई शताब्दी।
\textbf{7.} देशज--देशज युग्म संगमों का $(60/70)^{2} = 0.73$ हैं और उनकी
संतानें $F \approx 0.03$ रखती हैं (जमा कोई छोटी वृद्धि); बाक़ी सबका
$F = 0$ है: माध्य $F \approx 0.73\times 0.035 = 0.026$, यानी एक ही पीढ़ी
में चौथाई की गिरावट।
\textbf{8.} $F$ किसी व्यष्टि के भीतर वंश से पहचान है, और कोई भी
बहिःप्रजनित संतान उसे शून्य पर लौटा देती है; विषमयुग्मजता
युग्मविकल्पी-बारंबारताओं का गुण है, और प्रवासियों के युग्मविकल्पी केवल
आगे बढ़ाए जाने के साथ ही फैलते हैं। अब भी छोटी किसी समष्टि में अपवाह
उन्हें दसियों पीढ़ियों के भीतर फिर हटा देगा, इसलिए बचाव दोहराना होगा या
समष्टि बढ़ानी होगी।
\textbf{9.} बराबर लिंगों के लिए $N_{e} = N$: 500 वयस्क। 70 से $r =
0.05$ पर: $\ln(500/70)/0.05 = 39$ वर्ष।
\textbf{10.} आदर्श समष्टि में प्वासों कुल-आकार होते हैं, जिनका प्रसरण दो
के माध्य के बराबर होता है; $N_{e} = 4N/(2 + V_{k})$, इसलिए हर युग्म को
बराबर योगदान दिलाना ($V_{k} = 0$) $N_{e} = 2N$ देता है: घोंसले के भाग्य
से कोई वंशक्रम खोता नहीं।
\textbf{11.} कृत्रिम वीर्यसेचन के लिए उसका वीर्य जमा कर रखिए, और उसे कई
वर्षों भर कई असंबंधित मादाओं पर काम में लीजिए --- लागत यह है कि उसके
युग्मविकल्पी किसी छोटे कोश पर हावी हो सकते हैं, इसलिए उसका हिस्सा
$1/N_{e}$ के पास सीमित रखना होगा; या बाद में क्लोनन या युग्मक-उत्पादन के
लिए ऊतक हिमपरिरक्षित कीजिए, जिसकी लागत देरी तथा तकनीकी जोखिम है। कुछ न
करना उस वंशक्रम के युग्मविकल्पी सदा के लिए खो देता है।
\textbf{12.} कोई 430 पक्षियों के लिए \qty{160}{} मिलियन डॉलर: यानी प्रति
पक्षी लगभग \num{370000}। प्रति तोता महँगा, किसी जाति के सामने सस्ता ---
और किसी मोटरमार्ग के कुछ किलोमीटर से कम, और बजट असल में यही तुलना करते
हैं।
\textbf{13.} $(300/2000)^{0.25} = 0.62$: 180 में से लगभग 112 जातियाँ रखी
गईं, यानी 68 का विलोपन-ऋण।
\textbf{14.} हर खंड $(100/2000)^{0.25}\times 180 = 85$ जातियाँ। पूरी तरह
भिन्न: 255 तक, जो 180 के कोश से ऊपर असंभव है; एक-सी: 85। जो तय करता है
वह यह है कि खंडों के आवास कितने भिन्न हैं और \qty{100}{km^2} में कितनी
जातियाँ बनी रह सकती हैं; वास्तविक रूप से 85 और 112 के बीच।
\textbf{15.} $(320/2000)^{0.25}\times 180 = 114$, जबकि तीन अलग-थलग एक-सी
जातियों वाले खंडों के लिए 85: वह गलियारा पुनर्बसाव तथा जीन-प्रवाह को तीन
छोटी समष्टियों को किसी एक बड़ी समष्टि जैसा बरतवाने देता है।
\textbf{16.} बराबर लिंगों के साथ $N_{e} = 50$ यानी 25 युग्म:
\qty{1250}{km^2}। पूरा \qty{300}{km^2} छह युग्म थामता है,
$N_{e} \approx 12$; कोई खंड, दो। दोनों में से कोई पर्याप्त नहीं: जातियों
की गिनती चाहे जो हो, वह परभक्षी खो जाता है, और उसे जितना क्षेत्रफल चाहिए
उसकी रक्षा करना उससे छोटी हर चीज़ की रक्षा करता है --- यानी कोई छत्र
जाति।
\textbf{17.} $\ln(200/10)/0.03 = 100$ वर्ष: वह ऋण किसी शताब्दी भर चुकाया
जाता है, वन कटने के बहुत बाद तक।
\textbf{18.} किनारे हर खंड से भीतरी आवास छीन लेते हैं, इसलिए काम आने
लायक़ क्षेत्रफल मानचित्रित क्षेत्रफल से कम होता है और हानि उस नियम की
भविष्यवाणी से अधिक; दूसरी बार उगी वनस्पति या बाड़-पंक्तियों का कोई
आधात्री, जिसे कुछ जातियाँ काम में ले सकें, प्रभावी क्षेत्रफल जोड़ता है
और हानि को छोटा कर देता है।
\textbf{19.} $r = \ln(100/31)/8 = \qty{0.146}{yr^{-1}}$; दुगुना होने का
समय $\ln 2/0.146 = 4.7$ वर्ष।
\textbf{20.} $N(8) = 120/[1 + (120/31 - 1)\mathrm{e}^{-1.17}] = 120/(1 +
2.87\times 0.31) = 63$। देखा गया 100 वही चरघातांकी मान है, क्योंकि
वृद्धि-दर इन्हीं दो बिंदुओं पर फिट की गई थी, न कि संभार्य 63: पहले
वर्षों में भेड़ियों को कोई घनत्व-निर्भरता मिली ही नहीं --- ख़ाली क्षेत्र
तथा प्रचुर एल्क।
\textbf{21.} $\ln(4000/\num{17000})/20 = -0.072$: यानी प्रति वर्ष
\qty{7}{\%}। सौ भेड़िये जो \num{8000} के औसत वाले किसी झुंड से प्रति वर्ष
\num{2000} एल्क लेते हों, वह प्रति वर्ष चौथाई है --- उस ह्रास को अकेले
चलाने के लिए पर्याप्त से अधिक, यद्यपि उद्यान के बाहर का शिकार, सूखा तथा
बछड़ों पर ग्रिज़ली का परभक्षण उसमें साझीदार रहे।
\textbf{22.} साठ परिपक्व व्यष्टि, 250 से नीचे: अतिसंकटग्रस्त।
$N_{e} \approx 12$ पर वह परभक्षी (कुछ दर्जन प्राणी): वह भी
अतिसंकटग्रस्त।
\textbf{23.} 44 वर्षों में $3\times 10^{9}$ से $1$ तक: $r = -\ln(3\times
10^{9})/44 = -0.50$ प्रति वर्ष, यानी हर सत्रह महीने में आधा। कोई कटाई
हर वर्ष अरबों का आधा नहीं हटाती; वह गिरावट पहले धीमी थी और अंत में कहीं
अधिक तेज़, क्योंकि कोई बस्ती बनाकर प्रजनन करने वाला किसी देहली घनत्व से
नीचे जनन करना बंद कर देता है जबकि उसके परभक्षी अब अभिभूत नहीं होते ---
\omterm{def:b3:conservation-biology:small}{ऐली प्रभाव} ने किसी ह्रास को किसी ढहाव में बदल दिया।
\textbf{24.} 1850 में वह समष्टि अरबों में थी, उसकी मापी गई वृद्धि-दर
धनात्मक और उसका प्रसरण छोटा, और कोई व्यवहार्यता विश्लेषण उसी गतिकी का
प्रक्षेप करता है जो उसे दी जाए: वह जोखिम को शून्य ठहराता। वह उस बाहरी
प्रणोदन से चूक जाता --- रेलमार्ग तथा तार, जिन्होंने शिकारियों को हर
बसेरे तक पीछा करने और उसे ख़ाली करने दिया, और वन-सफ़ाई --- और किसी बस्ती
बनाकर प्रजनन करने वाले में उस ऐली देहली से, जिनमें से कोई भी किसी प्रचुर
जाति की जनसांख्यिकी में नहीं होता; कोई PVA बहिर्वेशन करता है, और विलोपन
के कारण प्रायः नए होते हैं।
\textbf{25.} आज $N_{e} \approx 53$; 60 पर दस और पीढ़ियों बाद
$F \approx 0.12$ (सुयोग्यता का \qty{7}{\%}); वह सिकुड़ा हुआ वन अपनी जातियों
का लगभग \qty{62}{\%} रखता है।
\end{solution}
